\section{Injective resolutions}
\label{section-injective-resolutions}

\noindent
In this section we prove some lemmas regarding the existence
of injective resolutions in abelian categories having enough injectives.

\begin{definition}
\label{definition-injective-resolution}
Let $\mathcal{A}$ be an abelian category.
Let $A \in \Ob(\mathcal{A})$.
An {\it injective resolution of $A$} is a complex
$I^\bullet$ together with a map $A \to I^0$ such
that:
\begin{enumerate}
\item We have $I^n = 0$ for $n < 0$.
\item Each $I^n$ is an injective object of $\mathcal{A}$.
\item The map $A \to I^0$ is an isomorphism onto $\Ker(d^0)$.
\item We have $H^i(I^\bullet) = 0$ for $i > 0$.
\end{enumerate}
Hence $A[0] \to I^\bullet$ is a quasi-isomorphism.
In other words the complex
$$
\ldots \to 0 \to A \to I^0 \to I^1 \to \ldots
$$
is acyclic.
Let $K^\bullet$ be a complex in $\mathcal{A}$.
An {\it injective resolution of $K^\bullet$} is a complex
$I^\bullet$ together with a map $\alpha : K^\bullet \to I^\bullet$
of complexes such that
\begin{enumerate}
\item We have $I^n = 0$ for $n \ll 0$, i.e., $I^\bullet$ is bounded below.
\item Each $I^n$ is an injective object of $\mathcal{A}$.
\item The map $\alpha : K^\bullet \to I^\bullet$ is a
quasi-isomorphism.
\end{enumerate}
\end{definition}

\noindent
In other words an injective resolution $K^\bullet \to I^\bullet$
gives rise to a diagram
$$
\xymatrix{
\ldots \ar[r] & K^{n - 1} \ar[d] \ar[r] & K^n \ar[d] \ar[r] &
K^{n + 1} \ar[d] \ar[r] & \ldots \\
\ldots \ar[r] & I^{n - 1} \ar[r] & I^n \ar[r] & I^{n + 1} \ar[r] & \ldots
}
$$
which induces an isomorphism on cohomology objects in each degree.
An injective resolution of an object $A$ of $\mathcal{A}$
is almost the same thing as an injective resolution of
the complex $A[0]$.

\begin{lemma}
\label{lemma-cohomology-bounded-below}
Let $\mathcal{A}$ be an abelian category.
Let $K^\bullet$ be a complex of $\mathcal{A}$.
\begin{enumerate}
\item If $K^\bullet$ has an injective resolution then
$H^n(K^\bullet) = 0$ for $n \ll 0$.
\item If $H^n(K^\bullet) = 0$ for all $n \ll 0$ then there
exists a quasi-isomorphism $K^\bullet \to L^\bullet$
with $L^\bullet$ bounded below.
\end{enumerate}
\end{lemma}

\begin{proof}
Omitted. For the second statement use
$L^\bullet = \tau_{\geq n}K^\bullet$ for
some $n \ll 0$. See
Homology, Section \ref{homology-section-truncations}
for the definition of the truncation $\tau_{\geq n}$.
\end{proof}

\begin{lemma}
\label{lemma-injective-resolutions-exist}
Let $\mathcal{A}$ be an abelian category.
Assume $\mathcal{A}$ has enough injectives.
\begin{enumerate}
\item Any object of $\mathcal{A}$ has an injective resolution.
\item If $H^n(K^\bullet) = 0$ for all $n \ll 0$ then
$K^\bullet$ has an injective resolution.
\item If $K^\bullet$ is a complex with $K^n = 0$ for $n < a$, then
there exists an injective resolution $\alpha : K^\bullet \to I^\bullet$
with $I^n = 0$ for $n < a$ such that each $\alpha^n : K^n \to I^n$ is
injective.
\end{enumerate}
\end{lemma}

\begin{proof}
Proof of (1). First choose an injection $A \to I^0$ of $A$ into an
injective object of $\mathcal{A}$. Next, choose an injection
$I_0/A \to I^1$ into an injective object of $\mathcal{A}$.
Denote $d^0$ the induced map $I^0 \to I^1$.
Next, choose an injection $I^1/\Im(d^0) \to I^2$ into
an injective object of $\mathcal{A}$. Denote $d^1$ the induced
map $I^1 \to I^2$. And so on.
By Lemma \ref{lemma-cohomology-bounded-below} part (2) follows from part (3).
Part (3) is a special case of
Lemma \ref{lemma-subcategory-right-resolution}.
\end{proof}

\begin{lemma}
\label{lemma-acyclic-is-zero}
Let $\mathcal{A}$ be an abelian category.
Let $K^\bullet$ be an acyclic complex.
Let $I^\bullet$ be bounded below and consisting of injective objects.
Any morphism $K^\bullet \to I^\bullet$ is homotopic to zero.
\end{lemma}

\begin{proof}
Let $\alpha : K^\bullet \to I^\bullet$ be a morphism of
complexes. Note that $\alpha^j = 0$ for $j \ll 0$ as $I^\bullet$
is bounded below. In particular, we can find an $n$ such that
there exist $h^j :  K^j \to I^{j - 1}$ for $j \leq n$ such that
$\alpha^j = d^{j - 1} \circ h^j + h^{j + 1} \circ d^j$
for $j < n$. We will show that there exists a morphism
$h^{n + 1} : K^{n + 1} \to I^n$
such that $\alpha^n = d^{n - 1} \circ h^n + h^{n + 1} \circ d^n$.
Note that
\begin{align*}
(\alpha^n - d^{n - 1} \circ h^{n - 1}) \circ d^{n - 1}
& =
\alpha^{n - 1} \circ d^{n - 1} - d^{n - 1} \circ h^{n - 1} \circ d^{n - 1} \\
& =
d^{n - 1} \circ \alpha^{n - 1} - d^{n - 1} \circ h^{n - 1} \circ d^{n - 1} \\
& =
d^{n - 1} \circ (d^{n - 2} \circ h^{n - 1} + h^n \circ d^{n - 1})
- d^{n - 1} \circ h^{n - 1} \circ d^{n - 1} \\
& = 0
\end{align*}
Since $K^\bullet$ is acyclic we have
$d^{n - 1}(K^{n - 1}) = \Ker(K^n \to K^{n + 1})$.
Hence we can think of $\alpha^n -  d^{n - 1} \circ h^{n - 1}$
as a map into $I^n$ defined on the subobject $\Im(K^n \to K^{n + 1})$
of $K^{n + 1}$. By injectivity of the object $I^n$ we can extend this to
a map $h^{n + 1} : K^{n + 1} \to I^n$. With this choice the reader
checks that we indeed have
$\alpha^n = d^{n - 1} \circ h^n + h^{n + 1} \circ d^n$.

\medskip\noindent
By induction on $n$ we conclude we can find $h  = (h^j)_{j \in \mathbf{Z}}$
which forms a homotopy between $\alpha$ and $0$ as desired.
\end{proof}

\begin{remark}
\label{remark-easier-proofs}
Let $\mathcal{A}$ be an abelian category.
Using the fact that $K(\mathcal{A})$ is a triangulated category we
may use
Lemma \ref{lemma-acyclic-is-zero}
to obtain proofs of some of the lemmas below which are usually proved by
chasing through diagrams.
Namely, suppose that $\alpha : K^\bullet \to L^\bullet$ is a quasi-isomorphism
of complexes. Then
$$
(K^\bullet, L^\bullet, C(\alpha)^\bullet, \alpha, i, -p)
$$
is a distinguished triangle in $K(\mathcal{A})$
(Lemma \ref{lemma-the-same-up-to-isomorphisms})
and $C(\alpha)^\bullet$ is an acyclic complex
(Lemma \ref{lemma-acyclic}).
Next, let $I^\bullet$ be a bounded below complex of injective objects. Then
$$
\xymatrix{
\Hom_{K(\mathcal{A})}(C(\alpha)^\bullet, I^\bullet) \ar[r] &
\Hom_{K(\mathcal{A})}(L^\bullet, I^\bullet) \ar[r] &
\Hom_{K(\mathcal{A})}(K^\bullet, I^\bullet) \ar[lld] \\
\Hom_{K(\mathcal{A})}(C(\alpha)^\bullet[-1], I^\bullet)
}
$$
is an exact sequence of abelian groups, see
Lemma \ref{lemma-representable-homological}.
At this point
Lemma \ref{lemma-acyclic-is-zero}
guarantees that the outer two groups are zero and hence
$\Hom_{K(\mathcal{A})}(L^\bullet, I^\bullet) =
\Hom_{K(\mathcal{A})}(K^\bullet, I^\bullet)$.
\end{remark}

\begin{lemma}
\label{lemma-morphisms-lift}
Let $\mathcal{A}$ be an abelian category.
Consider a solid diagram
$$
\xymatrix{
K^\bullet \ar[r]_\alpha \ar[d]_\gamma & L^\bullet \ar@{-->}[dl]^\beta \\
I^\bullet
}
$$
where $I^\bullet$ is bounded below and consists of injective
objects, and $\alpha$ is a quasi-isomorphism.
\begin{enumerate}
\item There exists a map of complexes $\beta$ making the diagram
commute up to homotopy.
\item If $\alpha$ is injective in every degree
then we can find a $\beta$ which makes the diagram commute.
\end{enumerate}
\end{lemma}

\begin{proof}
The ``correct'' proof of part (1) is explained in
Remark \ref{remark-easier-proofs}.
We also give a direct proof here.

\medskip\noindent
We first show that (2) implies (1). Namely, let
$\tilde \alpha : K \to \tilde L^\bullet$, $\pi$, $s$ be as in
Lemma \ref{lemma-make-injective}. Since $\tilde \alpha$ is injective
by (2) there exists a morphism $\tilde \beta : \tilde L^\bullet \to I^\bullet$
such that $\gamma = \tilde \beta \circ \tilde \alpha$. Set
$\beta = \tilde \beta \circ s$. Then we have
$$
\beta \circ \alpha
=
\tilde \beta \circ s \circ \pi \circ \tilde \alpha
\sim
\tilde \beta \circ \tilde \alpha
=
\gamma
$$
as desired.

\medskip\noindent
Assume that $\alpha : K^\bullet \to L^\bullet$ is injective.
Suppose we have already defined $\beta$ in all degrees
$\leq n - 1$ compatible with differentials and such that
$\gamma^j = \beta^j \circ \alpha^j$ for all $j \leq n - 1$.
Consider the commutative solid diagram
$$
\xymatrix{
K^{n - 1} \ar[r] \ar@/_2pc/[dd]_\gamma \ar[d]^\alpha &
K^n \ar@/^2pc/[dd]^\gamma \ar[d]^\alpha \\
L^{n - 1} \ar[r] \ar[d]^\beta &
L^n \ar@{-->}[d] \\
I^{n - 1} \ar[r] &
I^n
}
$$
Thus we see that the dotted arrow is prescribed on the subobjects
$\alpha(K^n)$ and $d^{n - 1}(L^{n - 1})$. Moreover, these two arrows
agree on $\alpha(d^{n - 1}(K^{n - 1}))$. Hence if
\begin{equation}
\label{equation-qis}
\alpha(d^{n - 1}(K^{n - 1}))
=
\alpha(K^n) \cap d^{n - 1}(L^{n - 1})
\end{equation}
then these morphisms glue to a morphism
$\alpha(K^n) + d^{n - 1}(L^{n - 1}) \to I^n$ and, using the injectivity
of $I^n$, we can extend this to a morphism from all of $L^n$ into $I^n$.
After this by induction we get the morphism $\beta$ for all $n$ simultaneously
(note that we can set $\beta^n = 0$ for all $n \ll 0$ since $I^\bullet$
is bounded below -- in this way starting the induction).

\medskip\noindent
It remains to prove the equality (\ref{equation-qis}).
The reader is encouraged to argue this for themselves with a suitable
diagram chase. Nonetheless here is our argument.
Note that the inclusion
$\alpha(d^{n - 1}(K^{n - 1})) \subset \alpha(K^n) \cap d^{n - 1}(L^{n - 1})$
is obvious. Take an object $T$ of $\mathcal{A}$ and a morphism
$x : T \to L^n$ whose image is contained in the subobject
$\alpha(K^n) \cap d^{n - 1}(L^{n - 1})$.
Since $\alpha$ is injective we see that $x = \alpha \circ x'$ for
some $x' : T \to K^n$. Moreover, since $x$ lies in $d^{n - 1}(L^{n - 1})$
we see that $d^n \circ x = 0$. Hence using injectivity of $\alpha$ again
we see that $d^n \circ x' = 0$. Thus $x'$ gives a morphism
$[x'] : T \to H^n(K^\bullet)$. On the other hand the corresponding
map $[x] : T \to H^n(L^\bullet)$ induced by $x$ is zero by assumption.
Since $\alpha$ is a quasi-isomorphism we conclude that $[x'] = 0$.
This of course means exactly that the image of $x'$ is
contained in $d^{n - 1}(K^{n - 1})$ and we win.
\end{proof}

\begin{lemma}
\label{lemma-morphisms-equal-up-to-homotopy}
Let $\mathcal{A}$ be an abelian category.
Consider a solid diagram
$$
\xymatrix{
K^\bullet \ar[r]_\alpha \ar[d]_\gamma & L^\bullet \ar@{-->}[dl]^{\beta_i} \\
I^\bullet
}
$$
where $I^\bullet$ is bounded below and consists of injective
objects, and $\alpha$ is a quasi-isomorphism.
Any two morphisms $\beta_1, \beta_2$ making the diagram commute
up to homotopy are homotopic.
\end{lemma}

\begin{proof}
This follows from
Remark \ref{remark-easier-proofs}.
We also give a direct argument here.

\medskip\noindent
Let $\tilde \alpha : K \to \tilde L^\bullet$, $\pi$, $s$ be as in
Lemma \ref{lemma-make-injective}. If we can show that $\beta_1 \circ\pi$
is homotopic to $\beta_2 \circ \pi$, then we deduce that
$\beta_1 \sim \beta_2$ because $\pi \circ s$ is the identity.
Hence we may assume $\alpha^n : K^n \to L^n$ is the
inclusion of a direct summand for all $n$. Thus we get a
short exact sequence of complexes
$$
0 \to K^\bullet \to L^\bullet \to M^\bullet \to 0
$$
which is termwise split and such that $M^\bullet$ is acyclic.
We choose splittings $L^n = K^n \oplus M^n$, so we have
$\beta_i^n : K^n \oplus M^n \to I^n$ and $\gamma^n : K^n \to I^n$.
In this case the condition on $\beta_i$ is that there are morphisms
$h_i^n : K^n \to I^{n - 1}$ such that
$$
\gamma^n - \beta_i^n|_{K^n} = d \circ h_i^n + h_i^{n + 1} \circ d
$$
Thus we see that
$$
\beta_1^n|_{K^n} - \beta_2^n|_{K^n}
=
d \circ (h_1^n - h_2^n) + (h_1^{n + 1} - h_2^{n + 1}) \circ d
$$
Consider the map $h^n : K^n \oplus M^n \to I^{n - 1}$ which
equals $h_1^n - h_2^n$ on the first summand and zero on the second.
Then we see that
$$
\beta_1^n - \beta_2^n
-
(d \circ h^n + h^{n + 1}) \circ d
$$
is a morphism of complexes $L^\bullet \to I^\bullet$
which is identically zero on the subcomplex $K^\bullet$.
Hence it factors as $L^\bullet \to M^\bullet \to I^\bullet$.
Thus the result of the lemma follows from Lemma \ref{lemma-acyclic-is-zero}.
\end{proof}

\begin{lemma}
\label{lemma-morphisms-into-injective-complex}
Let $\mathcal{A}$ be an abelian category.
Let $I^\bullet$ be bounded below complex consisting of injective
objects. Let $L^\bullet \in K(\mathcal{A})$. Then
$$
\Mor_{K(\mathcal{A})}(L^\bullet, I^\bullet)
=
\Mor_{D(\mathcal{A})}(L^\bullet, I^\bullet).
$$
\end{lemma}

\begin{proof}
Let $a$ be an element of the right hand side.
We may represent $a = \gamma\alpha^{-1}$ where
$\alpha : K^\bullet \to L^\bullet$
is a quasi-isomorphism and $\gamma : K^\bullet \to I^\bullet$ is a map
of complexes. By
Lemma \ref{lemma-morphisms-lift}
we can find a morphism $\beta : L^\bullet \to I^\bullet$ such that
$\beta \circ \alpha$ is homotopic to $\gamma$. This proves that the
map is surjective. Let $b$ be an element of the left hand side
which maps to zero in the right hand side. Then $b$ is the homotopy class
of a morphism $\beta : L^\bullet \to I^\bullet$ such that there exists a
quasi-isomorphism $\alpha : K^\bullet \to L^\bullet$ with
$\beta  \circ \alpha$ homotopic to zero. Then
Lemma \ref{lemma-morphisms-equal-up-to-homotopy}
shows that $\beta$ is homotopic to zero also, i.e., $b = 0$.
\end{proof}

\begin{lemma}
\label{lemma-injective-resolution-ses}
Let $\mathcal{A}$ be an abelian category.
Assume $\mathcal{A}$ has enough injectives.
For any short exact sequence
$0 \to A^\bullet \to B^\bullet \to C^\bullet \to 0$
of $\text{Comp}^{+}(\mathcal{A})$ there exists a
commutative diagram in $\text{Comp}^{+}(\mathcal{A})$
$$
\xymatrix{
0 \ar[r] &
A^\bullet \ar[r] \ar[d] &
B^\bullet \ar[r] \ar[d] &
C^\bullet \ar[r] \ar[d] &
0 \\
0 \ar[r] &
I_1^\bullet \ar[r] &
I_2^\bullet \ar[r] &
I_3^\bullet \ar[r] &
0
}
$$
where the vertical arrows are injective resolutions and
the rows are short exact sequences of complexes. Additionally,
\begin{enumerate}
\item given any injective resolution $A^\bullet \to I^\bullet$
we may assume $I_1^\bullet = I^\bullet$,
\item if $A^n = B^n = C^n = 0$ for $n < 0$, then we may
assume $I_j^n = 0$ for $n < 0$,
\item we can combine (1) and (2) if also $I^n = 0$ for $n < 0$, and
\item add more here.
\end{enumerate}
\end{lemma}

\begin{proof}
Step 1. Choose an injective resolution $A^\bullet \to I^\bullet$ (see
Lemma \ref{lemma-injective-resolutions-exist}) or use the given one.
Recall that $\text{Comp}^{+}(\mathcal{A})$ is an
abelian category, see
Homology, Lemma \ref{homology-lemma-cat-cochain-abelian}.
Hence we may form the pushout along
the map $A^\bullet \to I^\bullet$ to get
$$
\xymatrix{
0 \ar[r] &
A^\bullet \ar[r] \ar[d] &
B^\bullet \ar[r] \ar[d] &
C^\bullet \ar[r] \ar[d] &
0 \\
0 \ar[r] &
I^\bullet \ar[r] &
E^\bullet \ar[r] &
C^\bullet \ar[r] &
0
}
$$
Because of the $5$-lemma and the last assertion of
Homology, Lemma \ref{homology-lemma-long-exact-sequence-cochain}
the map $B^\bullet \to E^\bullet$ is a quasi-isomorphism.
Note that the lower short exact sequence is termwise split, see
Homology, Lemma \ref{homology-lemma-characterize-injectives}.
Hence it suffices to prove the lemma when
$0 \to A^\bullet \to B^\bullet \to C^\bullet \to 0$ is
termwise split.

\medskip\noindent
Step 2. Choose splittings. In other words, write
$B^n = A^n \oplus C^n$. Denote $\delta : C^\bullet \to A^\bullet[1]$
the morphism as in
Homology, Lemma \ref{homology-lemma-ses-termwise-split-cochain}.
Choose injective resolutions $f_1 : A^\bullet \to I_1^\bullet$
and $f_3 : C^\bullet \to I_3^\bullet$. (If $A^\bullet$ is a complex of
injectives, then use $I_1^\bullet = A^\bullet$. If $A^n = C^n = 0$ for
$n < 0$, then pick these such that $I_1^n = I_3^n = 0$ for $n < 0$, see
Lemma \ref{lemma-injective-resolutions-exist}.)
We may assume $f_3$ is injective in
every degree. By Lemma \ref{lemma-morphisms-lift} we may find
a morphism $\delta' : I_3^\bullet \to I_1^\bullet[1]$ such
that $\delta' \circ f_3 = f_1[1] \circ \delta$ (equality of
morphisms of complexes). Set $I_2^n = I_1^n \oplus I_3^n$.
Define
$$
d_{I_2}^n =
\left(
\begin{matrix}
d_{I_1}^n & (\delta')^n \\
0 & d_{I_3}^n
\end{matrix}
\right)
$$
and define the maps $B^n \to I_2^n$ to be given as the
sum of the maps $A^n \to I_1^n$ and $C^n \to I_3^n$.
Everything is clear.
\end{proof}








